\documentclass[10pt,a4paper]{amsart}
\usepackage[margin=19mm]{geometry}
\usepackage{amsmath,amssymb,mathtools}
\usepackage[scr=rsfs,cal=euler]{mathalfa}
\usepackage{xfrac}
\usepackage[unicode,hidelinks,bookmarks=false]{hyperref}
\newcommand{\ie}{i.e.\ }
\newcommand{\cf}{cf.\ }
\newcommand{\resp}{respectively\ }
\newcommand{\Subsection}{\subsection}
\providecommand{\nobreakdash}{\nobreak\hbox{-}}
\setlength{\emergencystretch}{2em}
\multlinegap0pt
\usepackage{amssymb}
\newtheorem{proposition}{Proposition}
\newtheorem{theorem}{Theorem}
\newcommand {\bC} {\mathbb {C}}
\newcommand {\bP} {\mathbb {P}}
\newcommand {\bR} {\mathbb R}
\newcommand {\cP} {\mathcal P}
\title{Canonical forms and moment-generating functions of plane polypols}
\author{Boris Shapiro}
\date{Source: arXiv:2605.10864v2 (CC BY 4.0)}
\begin{document}
\maketitle

\noindent\emph{Source and licence.} Canonical forms and moment-generating functions of plane polypols. Version arXiv:2605.10864v2 (CC BY 4.0), distributed by arXiv under the Creative Commons Attribution 4.0 International licence, http://creativecommons.org/licenses/by/4.0/, as recorded in the arXiv licence metadata for this exact version and shown on the abstract page https://arxiv.org/abs/2605.10864v2. Full licence notice: \texttt{licenses/arxiv-2605.10864-CC-BY-4.0.txt}.\par\smallskip

\noindent\textit{Editorial scope.} Counted unit: the article's theorem th:GenFucnt (source0.tex lines 267--273). The article's complete original proof of it is delivered with it (source0.tex lines 317--392, one \texttt{proof} environment of 76 lines, which the article places away from the statement). No mathematics is written, repaired or re-derived: the statement and the proof are the article's own lines, in the article's own notation, and no retained line is substituted or re-flowed.  Retained uncounted as the source-local context the proof needs: lines 239--246.  Nothing was omitted from any retained span, so the package carries no omission record.  No retained line contains a citation, so the wrapper carries no bibliography and no bibliography entry.  No retained line contains a figure environment, an image-inclusion command or drawing code, so no image is delivered and the package declares no asset dependency.  Editorial formatting only, in addition: an AMS \texttt{amsart} preamble that restates the article's own declarations and macros used by the retained lines, namely the environments \texttt{proposition}, \texttt{theorem} on separate counters, together with the standard packages those lines need.  The retained environments print in this document as Proposition 1, Theorem 1 in order of appearance, whereas the article prints the same items with the numbering of its own full document.  The complete native source of the article as distributed in the arXiv e-print for this exact version is bundled unchanged as \texttt{sources/200326/source0.tex}.
\begin{proposition}[Polytope case]\label{prop:Poly}
Let $P\subset \bR^N$ be a full-dimensional polytope with vertices $V(P)$, and assume that the origin is in the interior of $P$ so that the polar body is bounded.  Then
$$
F_P(u)=\frac{A_{P^\circ}(u)}{\prod_{v\in V(P)}(1-\langle v,u\rangle)},
$$
where $A_{P^\circ}$ is the adjoint numerator of the polar polytope.  Equivalently,
$F_P(u)du_1\wedge\cdots\wedge du_N$ is the canonical form of $P^\circ$, up to the normalization fixed by orientation.
\end{proposition}

\begin{theorem}[Fantappi\`e transform of a rational polypol]\label{th:GenFucnt}
Let $\cP\subset\bR^2$ be a regular rational polypol and set
$$
F_{\cP}(u,v)=2\int_{\cP}\frac{dxdy}{(1-ux-vy)^3}.
$$
Then $F_{\cP}$ is obtained by summing explicit one-dimensional integrals over rational boundary arcs.  Consequently it is an elementary multivalued function of $(u,v)$, built from algebraic functions, rational functions, logarithms and arctangent terms; equivalently it is a period of the rational differential form above and is holonomic in the parameters.  If $\cP$ is a polygon, all boundary integrals are rational and Proposition~\ref{prop:Poly} identifies $F_{\cP}$ with the canonical form of the polar polygon.  For genuinely curved boundary arcs one should not expect rationality; the upper half-disk example below already contains algebraic and arctangent terms, while the full disk gives the especially simple smooth-boundary model $2\pi(1-u^2-v^2)^{-3/2}$.
\end{theorem}

\begin{proof}[Proof of Theorem~\ref{th:GenFucnt}]
Put
\[
        Q_{u,v}(x,y)=1-ux-vy.
\]
The formula is first proved for $(u,v)$ in a sufficiently small neighbourhood of the origin.  In this region $Q_{u,v}$ has no zero on $\overline{\cP}$, and the integral defining $F_{\cP}$ is an ordinary convergent integral depending holomorphically on $(u,v)$.  The final statement then follows on the universal cover of the complement of the singular locus by analytic continuation.

Assume first that $u\ne 0$.  A direct calculation gives
\[
   d\left(\frac{dy}{2uQ_{u,v}(x,y)^2}\right)
       =\frac{dx\wedge dy}{Q_{u,v}(x,y)^3}.
\]
Indeed,
\[
 d\left(Q_{u,v}^{-2}dy\right)
 =\frac{2u\,dx\wedge dy}{Q_{u,v}^{3}}.
\]
Stokes' theorem therefore gives
\[
  \int_{\cP}\frac{dxdy}{Q_{u,v}(x,y)^3}
   =\frac{1}{2u}\int_{\partial\cP}\frac{dy}{Q_{u,v}(x,y)^2}.
\]
Equivalently, if $v\ne0$, using
\[
   d\left(-\frac{dx}{2vQ_{u,v}(x,y)^2}\right)
       =\frac{dx\wedge dy}{Q_{u,v}(x,y)^3},
\]
one obtains
\[
  \int_{\cP}\frac{dxdy}{Q_{u,v}(x,y)^3}
   =-\frac{1}{2v}\int_{\partial\cP}\frac{dx}{Q_{u,v}(x,y)^2}.
\]
Multiplying by $2$, according to the normalization in the definition of $F_{\cP}$, gives the boundary representations
\begin{equation}\label{eq:boundary-F-proof}
 F_{\cP}(u,v)=\frac{1}{u}\int_{\partial\cP}\frac{dy}{Q_{u,v}(x,y)^2}
              =-\frac{1}{v}\int_{\partial\cP}\frac{dx}{Q_{u,v}(x,y)^2},
\end{equation}
where the first identity is used on $u\ne0$ and the second on $v\ne0$.  Since both sides are holomorphic near the origin after taking the original area integral as the definition, these formulas determine the same analytic germ and extend across $u=0$ or $v=0$ by continuation.

Now decompose the oriented boundary into finitely many rational arcs
\[
        \partial\cP=\gamma_1\cup\cdots\cup\gamma_s.
\]
For each arc choose a rational parametrization
\[
        \gamma_j:\tau\in [a_j,b_j]\longmapsto (x_j(\tau),y_j(\tau)),
        \qquad x_j(\tau),y_j(\tau)\in\bR(\tau),
\]
compatible with the boundary orientation.  Substituting this parametrization in the first formula in \eqref{eq:boundary-F-proof} gives
\begin{equation}\label{eq:rational-arc-integrals}
 F_{\cP}(u,v)=\frac{1}{u}
       \sum_{j=1}^s\int_{a_j}^{b_j}
       \frac{y_j'(\tau)\,d\tau}
       {\bigl(1-u x_j(\tau)-v y_j(\tau)\bigr)^2},
\end{equation}
whenever $u\ne0$.  Similarly, for $v\ne0$,
\begin{equation}\label{eq:rational-arc-integrals-v}
 F_{\cP}(u,v)=-\frac{1}{v}
       \sum_{j=1}^s\int_{a_j}^{b_j}
       \frac{x_j'(\tau)\,d\tau}
       {\bigl(1-u x_j(\tau)-v y_j(\tau)\bigr)^2}.
\end{equation}
The integrands in \eqref{eq:rational-arc-integrals} and \eqref{eq:rational-arc-integrals-v} are rational functions of $\tau$ whose coefficients are rational functions of $(u,v)$.  Hence each boundary contribution is an explicitly computable integral of a rational differential on $\bP^1$.

By partial fractions over the algebraic closure of the coefficient field $\bC(u,v)$, an antiderivative of such a rational differential is a sum of a rational function of $\tau$ and logarithms of algebraic linear factors in $\tau$.  Over the real coefficient field, conjugate pairs of logarithms may equivalently be written as logarithmic and arctangent terms, with algebraic functions of $(u,v)$ appearing through the roots and discriminants of the denominator.  Evaluation at the endpoints $a_j,b_j$ therefore expresses $F_{\cP}(u,v)$, after analytic continuation, as an elementary multivalued function built from algebraic functions, rational functions, logarithms and arctangent terms.  This proves the asserted boundary-integral representation and the elementary character of the transform.

The holonomicity statement follows from the same representation.  Each summand in \eqref{eq:rational-arc-integrals} is a definite integral, over a fixed interval in the parameter $\tau$, of a rational function in the variables $(\tau,u,v)$.  Such parameter-dependent rational periods are holonomic functions of the parameters; equivalently, differentiating with respect to $u$ and $v$ produces a finite-dimensional module modulo exact rational differentials in $\tau$, which gives linear differential equations with polynomial coefficients annihilating the integral.  Summing over the finitely many arcs preserves holonomicity.

If $\cP$ is a polygon, each arc is a line segment.  On a segment $(x,y)=p+t(q-p)$, $0\le t\le1$, the denominator $Q_{u,v}(x,y)$ is affine in $t$, and the boundary integrals above are rational functions of $(u,v)$.  Summing over the sides gives a rational function.  Proposition~\ref{prop:Poly} identifies precisely this rational function, multiplied by $du\wedge dv$, with the canonical form of the polar polygon, up to the orientation convention fixed there.

Finally, curved rational arcs generally produce non-rational contributions because the rational denominator in the boundary parameter may have quadratic or higher degree factors.  As a smooth rational-boundary model, the full disk gives
\[
        F_{\mathrm{disk}}(u,v)=\frac{2\pi}{(1-u^2-v^2)^{3/2}},
\]
which is algebraic but not rational.  This shows that rationality is special to the polygonal case and completes the proof.
\end{proof}

\end{document}
